\sscp{\tsc{Sula-Buru} and \tsc{Ambon-Seram}}{07-03-02}
In lexically specific ways,
all languages in both \tsc{Sula-Buru}
and \tsc{Ambon-Seram} reflect PMP *j > **l in *pajay ‘rice’,
and *pija ‘how many?’ (see \srf{07-04-01-02},
and \srf{10-01} for conditioning environments for reflexes of
**l in specific languages).
All languages in both subgroups for which we have cognates reflect
PMP *j > {\0} in *qaləjaw ‘sun, day’.
Both of these lexically specific innovations are also found in
Banda and partially in the \tsc{East Seram} branch of \tsc{Seram-Tanimbar-Bomberai}.
In \tsc{Taliabu} *pija also reflects \it{l} rather than the expected
\it{y},\footnote{
		Data from Taliabu languages are sparse.
		None appear to have reflexes of *pajay or *qaləjaw.}
so these lexically specific innovations are not exclusive to \tsc{Sula-Buru}
and \tsc{Ambon-Seram}.
These are shown in \trf{07-05}.
Data from \tsc{Seram-Tanimbar-Bomberai} languages (\crf{ch:STB})
are included in \trf{07-05}
and \trf{07-06} to show the same sound changes or lexical patterns
are found, but not in all branches.

\begin{table}[p]
	\centering\tcp{Lexically specific split of PMP *j in central Maluku languages}{07-05}
	\begin{adjustbox}{totalheight=\textheight}
	\st{6pt}\begin{tabular}{llll}\lsptoprule
local gl.	&	rice	&	how many?	&	sun, day\\
PMP	&	*pa\tbr{j}ay	&	*pi\tbr{j}a	&	*qalə\tbr{j}aw\\
%\midrule
%\mc{4}{c}{\tsc{Taliabu}}\\
\midrule
Soboyo	&	\cg{(bira)}	&	{\hr}hi\tbr{l}a	&	\cg{(dina)}\\
Taliabu	&	\cg{(bira)}	&	{\hr}hi\tbr{l}a	&	\cg{(dina)}\\
Kadai	&	\cg{(bira)}	&		&	\cg{(dina)}\\
%\midrule
%\mc{4}{c}{\tsc{Sula-Buru}}\\
\midrule
Mangole	&	\cg{(pamasi)}	&	{\hr}pi\tbr{l}a	&	{\hr}lia\\
Sanana	&	\cg{(bira)}	&	{\hr}bet pi\tbr{l}a	&	{\hr}lea\\
Lisela	&	{\hr}ha\tbr{l}a (<Kayeli)	&	{\hr}pi\tbr{l}a-ni	&	\cg{(haŋa {\tl} haŋa-t)}\\
Buru	&	{\hr}pa\tbr{l}a	&	{\hr}pi\tbr{l}a	&	{\hr}lea\\
Ambelau	&	{\hr}fa\tbr{l}a	&	{\hr}fi\tbr{l}a {\tl} fia	&	{\hr}lea {\tl} leá-i\\
%\midrule
%\mc{4}{c}{\tsc{Ambon-Seram}}\\
\midrule
Kayeli	&	{\hr}ha\tbr{l}a	&	{\hr}hi\tbr{l}a {\tl} i\tbr{l}a	&	{\hr}lea\\
Boano	&	\hp{*h}a\tbr{l}a	&	{\hr}i\tbr{n}a	&	\\
Luhu	&	\hp{*h}a\tbr{l}a-i	&	{\hr}pitui\tbr{l}a	&	{\hr}le-mata-ne\\
Larike	&	\hp{*h}a\tbr{l}a	&	{\hr}i\tbr{d}a	&	{\hr}dia-mata\\
Asilulu	&	\hp{*h}a\tbr{l}a	&	{\hr}i\tbr{l}a	&	{\hr}lia-mata\\
Sepa	&	{\hr}ha\tbr{l}a	&	{\hr}hi\tbr{l}a	&	{\hr}ria-mata-n\\
Paulohi	&	{\hr}fa\tbr{l}a	&	{\hr}fi\tbr{l}a	&	{\hr}ʀia-mata-i\\
Kamarian	&	{\hr}ha\tbr{l}a	&	{\hr}wi\tbr{r}a	&	{\hr}re-mata-i\\
Haruku	&	{\hr}ha\tbr{l}a	&	{\hr}wi\tbr{r}a	&	{\hr}lea-mata-i\\
N. Wemale	&	{\hr}ha\tbr{l}a	&	{\hr}hi\tbr{l}a	&	{\hr}lia-mata-i\\
N. Alune	&	\hp{*h}a\tbr{l}a	&	{\hr}i\tbr{l}a	&	{\hr}le-mata-i\\
Yalahatan	&	{\hr}ha\tbr{l}a	&	\cg{(ehuisa)}	&	{\hr}lia-mata-i\\
S. Nuaulu	&	\hp{*h}a\tbr{n}a	&	{\hr}i\tbr{n}a	&	\cg{(rani-e)}\\
Manusela	&	{\hr}ha\tbr{l}a	&	{\hr}hi\tbr{l}a	&	{\hr}lea\\
Liana-Seti	&	{\hr}ha\tbr{l}a	&	{\hr}en-hi\tbr{l}a	&	{\hr}lea\\ \midrule
Banda	&	\hp{*h}a\tbr{l}a	&	{\hr}i\tbr{l}o	&	{\hr}lea\\
%\midrule
%\mc{4}{c}{\tsc{Seram-Tanimbar-Bomberai}}\\
\midrule
Bobot	&		&	{\hr}fi\tbr{l}a	&	{\hr}lia matan\\
Masiwang	&	{\hr}fa\tbr{s}a	&	{\hr}hi\tbr{l}a	&	\cg{(roman)}\\
Geser	&	{\hr}fa\tbr{s}a	&	{\hr}fi\tbr{s}	&	{\hr}o\tbr{l}a\\
Gorom	&	{\hr}ha\tbr{s}a	&	{\hr}hi\tbr{s}	&	{\hr}o\tbr{l}a\\
Kowiai	&	{\hr}ɸa\tbr{s}a	&	{\hr}ɸi\tbr{s}	&	{\hr}o\tbr{r}a mataɸut\\
Yamdena	&	{\hr}fa\tbr{s}-e	&	{\hr}fi\tbr{r}	&	{\hr}le\tbr{r}-e\\
Kei	&	\cg{(kokat)}	&	{\hr}n-fi\tbr{r}	&	{\hr}le\tbr{r}\\
Onin	&	{\hr}pa\tbr{s}a	&	{\hr}fi\tbr{r}as	&	{\hr}re\tbr{r}a\\
		\lspbottomrule
	\end{tabular}
\end{adjustbox}
\end{table}

In the lexicon, all \tsc{Sula-Buru}
and \tsc{Ambon-Seram} languages appear to take their word for
‘hundred’ from PCM **utum,\footnote{
PAMS **utu-ni ‘hundred’ \citep{Stresemann1927};
PCM **utun ‘hundred’ \citep[77]{Collins1983}.}
rather than PMP *Ratus.
By contrast, all three \tsc{Taliabu} languages
and some \tsc{Seram-Tanimbar-Bomberai} languages reflect PMP *Ratus.
All \tsc{Sula-Buru} and \tsc{Ambon-Seram} languages for which
we have data reflect **nakaw ‘steal’; none directly reflect PMP *takaw.
All \tsc{Sula-Buru} and \tsc{Ambon-Seram} languages for which
we have data reflect PMP *kədəŋ ‘stand,
stretch’; none reflect *diRi ‘stand’ as a verb (All \tsc{Ambon-Seram}
languages for which data is available do reflect *hadiRi as a
noun meaning ‘[main] house post’. See \trf{10-05}).
Only **utum seems to be an innovation shared exclusively by \tsc{Sula-Buru} and \tsc{Ambon-Seram}.
But numbers are among the domains of language that are easily borrowed.
These developments in the lexicon are shown in \trf{07-06}.

\begin{table}[p]\centering
	\tcp{Shared lexicon in central Maluku languages}{07-06}
	\begin{adjustbox}{totalheight=\textheight}
	\begin{threeparttable}
	\begin{tabular}{llll}\lsptoprule
local gl.	&	  steal	&	  hundred	&	  stand\\
Proto-form	&	**nakaw	&	**utum &	*kədəŋ \\ \midrule
%\mc{4}{c}{\tsc{Taliabu}}\\ \midrule
Soboyo	&	paka-nako	&	\hr\tcg{(ratu)}	&	koho-in\\
Taliabu	&	paka-nako	&	\hr\tcg{(ratu)}	&	kohoŋ\\
%\midrule
%Kadai	&		&	\hr\tcg{(yatu, ratu)}	&	\\
%\mc{4}{c}{\tsc{Sula-Buru}}\\ \midrule
\midrule
Mangole	&	bi-naka	&	\hr\tcg{(saka)}	&	keli, geli\\
Sanana	&	bil-naka	&	ota (*uCu{\#} > \it{oCa})	&	gehi (*d > \it{h})\\
Lisela	&		&	utun	&	kere\\
Buru	&	eʔnaka	&	utun	&	kere-k\\
Ambelau	&		&	urun-e (*t > \it{r})	&	\tcg{(en-βati)}\\
%\midrule
%\mc{4}{c}{\tsc{Ambon-Seram}}\\
\midrule
Kayeli	&		&	utun-e	&	\tcg{(en-dobo)}\\
%Kelang	&		&	utun-e	&	\\
Boano	&		&	utun-e	&	ele\\
Piru	&	\hr\tcg{(mamposi)}	&	utum	&	kele\\
Larike	&	pa-naʔa	&	hutun-a	&	kele\\
Asilulu	&	pa-naʔa	&	utun	&	kele\\
Paulohi	&		&	utun-i	&	ele\\
Kaibobo	&	a-naʔa, pa-naʔa	&	utum	&	olo\\
Kamarian	&	\hr\tcg{(‹paloe›)}	&	utum	&	ere\\
Tihulale	&	\hr\tcg{(amaloʔo)}	&	utum	&	ele\\
Rumakai	&	\hr\tcg{(amaloʔo)}	&	utum	&	ele\\
Waisamu	&	a-naʔa	&	utum	&	olo\\
Haruku	&	\hr\tcg{(maloʔo)}	&	utun ele	&	i-olo\\
Haruku (Karu)	&	\hr\tcg{(paloʔo)}	&	utum	&	olo\\
Saparua	&	mu-naʔa	&	utuɲo	&	olo\\
Nusalaut	&	mu-naʔa, puna	&	utuɲo	&	olo\\
Sepa	&		&	utun	&	olo\\
N. Wemale	&		&	utun-e	&	ele\\
N. Alune	&	\hr\tcg{(ndea)}	&	utun-e	&	kele\\
W. Alune	&	kamanaʔa	&	utum	&	kele\\
Yalahatan	&	ama-naʔa	&	hutun	&	ere\\
S. Nuaulu	&	kima-naka	&	utun	&	oo\\
Manusela	&	kamana	&	hutu-ni esa	&	oho\\
Liana-Seti	&	au-kama-na	&	utn es	&	au-kolo\\ \midrule
Banda	&		&	utuɲa	&	\tcg{mileri} 	\\
%\midrule
%\mc{4}{c}{\tsc{Seram-Tanimbar-Bomberai (East Seram)\su{†}}}\\
\midrule
Bobot\su{†}			&	ma-naa						&											&a-kolan\\
Masiwang	&	\hr\tcg{(kahleit, kahale)}	&	\tcg{(sati-a)}\su{‡}	&kolan\\
%Geser			&	\hr\tcg{(fakaleus)}				&	ratsa [raʧa]  	{\tcgr}	&ma-riri 	{\tcgr}	\\
%Gorom			&	\hr\tcg{(haleus)}					&	ratsa [raʧa], ratu	{\tcgr}	&ma-riri 	{\tcgr}	\\
%Watubela	&														&											{\tcgr}	&	-lili 	{\tcgr}	\\
%Kowiai		&	\hr\tcg{(na-ɸaraʔeus)}			&	ratsa, [raʧa] 	{\tcgr}	&	na-ma-rir 	{\tcgr}	\\
%Yamdena		&	\hr\tcg{(na-mnaŋ)}					&	raty						{\tcgr}	&	na-mdiry	{\tcgr}	\\
		\lspbottomrule
	\end{tabular}
		\begin{tablenotes}\tnsize
			\item[†] {Most \tsc{Seram-Tanimbar-Bomberai} languages
								for which data is available have reflexes
								of *Ratus for `hundred' and *(ma)-diRi for `stand'.}
			\item[‡] {Masiwang \it{sati-a} is regular from *Ratus with
								expected *R > \it{s} and final *u > \it{i}.}
			%\item[Δ]
			%\item[‖]
			%\item[¶]
		\end{tablenotes}
	\end{threeparttable}
\end{adjustbox}
\end{table}

While these six lexical innovations together are enough to suggest
Central Maluku (consisting of \tsc{Sula-Buru},
\tsc{Ambon-Seram}, Banda, and the western-most members of \tsc{Seram-Tanimbar-Bomberai})
could be considered a subgroup,
it is not a strong one.
Most of these innovations are not strictly exclusive to Central Maluku.
And as we show in \srf{07-04},
the patterned differences between \tsc{Sula-Buru}
and \tsc{Ambon-Seram} are significant,
suggesting that if they were ever together,
they split very early.

“PCM” (Proto-Central Maluku) reconstructions should be seen not
as something reflecting a parent language,
but rather forms that can account for data related by form
and meaning across a geographical region in which there has been
much movement and contact (see also \srf{02-02}).
Some PCM reconstructions are not valid for all subgroups or branches.
There are parallels, for example,
in forms that account for some words found in
both \tsc{Bima-Lembata} and \tsc{Timor-Babar} languages.
This issue is revisited in \srf{15-04-03}.

